\section{Starting point}
\label{sec:platform-start}

The only publicly available L2O benchmark is the code release that accompanies
the Primer of Chen et al.~\cite{chen2021primer}, Open-L2O. It is not a
platform in the sense of one program that trains and evaluates many methods;
it is a collection of the original authors' research codebases placed side by
side: six TensorFlow libraries, namely five model-free ones (L2O-DM together
with L2O-RNNProp, L2O-Scale, L2O-Entropy, L2O-Jacobian, L2O-Swarm) and one
model-based one (the LISTA family), plus a PyTorch library (L2O-Minimax) not
used in this thesis. Each has its own \texttt{train.py} and \texttt{evaluate.py}, its own
command-line conventions (two of the three families use \texttt{argparse}, the
third uses \texttt{absl.flags}), its own result files, and its own notion of
what a ``problem'' or a ``seed'' is. Several of the libraries also ship
top-level modules with the same names (\texttt{meta.py}, \texttt{util.py},
\texttt{problems.py}, \texttt{networks.py}), so two of them cannot even be
imported into the same Python process. Comparing two methods therefore meant
running two unrelated pipelines by hand and reconciling their outputs
afterwards, with nothing guaranteeing that both had seen the same problem
instances, the same starting points or the same evaluation protocol. As noted
in the pre-study (Section~\ref{sec:existing-benchmarks}), no alternative has
been published since.

The pre-study set four goals for extending this code into a platform
(Section~\ref{sec:planned-extensions}): (i) a unified pipeline through which
every method is trained and evaluated the same way; (ii) built-in
reproducibility and fairness; (iii) code that is atomised and reusable enough
for a newcomer to add a method or an experiment; and (iv) the data generator
and metric functions the experiments of this thesis require. This chapter
describes what was built for each. The framing matters: the platform
\emph{extends} Open-L2O rather than replacing it. The method implementations
that produce every number in Chapter~\ref{ch:exp2} are the Primer's own, run
unchanged inside a harness; what is new is everything around them.
Table~\ref{tab:platform-origin} separates the two.

% Table 3.1 -- what the platform inherited from Open-L2O and what this thesis added.
% Requires: booktabs, array. Repo paths are relative to the repository root.
\begin{table}[htbp]
\centering
\small
\caption{Components of the benchmarking platform, by origin. \emph{Inherited} components are the Primer's Open-L2O code release~\cite{chen2021primer}; \emph{migrated} means ported to TensorFlow~2 / Keras~3 without functional change (Appendix~\ref{app:migration}); \emph{added} components were written for this thesis.}
\label{tab:platform-origin}
\begin{tabular}{@{} p{0.30\linewidth} l p{0.44\linewidth} @{}}
\toprule
Component & Origin & Location / notes \\
\midrule
\multicolumn{3}{@{}l}{\emph{Method implementations}} \\
L2O-DM, L2O-RNNProp (one library) & inherited, migrated & \texttt{Model\_Free\_L2O/L2O-DM and L2O-RNNProp/} \\
L2O-Scale & inherited, migrated & \texttt{Model\_Free\_L2O/L2O-Scale/}; training and evaluation trees unified \\
L2O-Entropy, L2O-Jacobian & inherited, migrated & built on the unified L2O-Scale tree \\
L2O-Swarm & inherited, migrated & \texttt{Model\_Free\_L2O/L2O-Swarm/} \\
LISTA, ALISTA (model-based library) & inherited, migrated & \texttt{Model\_Base\_L2O/}; 7 further unrolled variants present but not wired \\
Classical optimizers (Adam, SGD, Adagrad) & inherited & run through a host library's evaluation harness \\
\addlinespace
\multicolumn{3}{@{}l}{\emph{Unified pipeline (goal i)}} \\
Method contract, registry, result containers, runner & added & \texttt{Benchmarking/core/} \\
Subprocess adapters, one per method family & added & \texttt{Benchmarking/adapters/} \\
Flag validation against the target script & added & \texttt{core/shell.py} \\
Command-line entry point and suites & added & \texttt{Benchmarking/run.py} \\
\addlinespace
\multicolumn{3}{@{}l}{\emph{Reproducibility and fairness (goal ii)}} \\
Seed threading into training and evaluation of every method & added & each library's \texttt{train.py}/\texttt{evaluate.py} accepts \texttt{--seed} \\
Deterministic full-pass evaluation over the test set & added & model-free \texttt{problems.py}, \texttt{evaluate\_*.py} \\
Shared per-instance starting point $x_0$ for every method & added & generator writes \texttt{*\_x0.npy}; LISTA/ALISTA layer~0 consumes it \\
Provenance capture (commands, git state, diffs, hashes, timings) & added & \texttt{core/runner.py}, \texttt{experiments/sweep.py} \\
Instances-fingerprint contract and determinism/fairness test suite & added & \texttt{core/fingerprint.py}, \texttt{Benchmarking/tests/} \\
Container image with pinned dependencies & added & \texttt{Dockerfile} \\
\addlinespace
\multicolumn{3}{@{}l}{\emph{Extensibility (goal iii)}} \\
Sweep driver (grid, resume, error isolation, summaries) & added & \texttt{experiments/sweep.py} \\
Runnable experiment template & added & \texttt{experiments/experiment\_template.py} \\
Contributor guide and reference documentation & added & \texttt{CONTRIBUTING.md}, \texttt{README.md} \\
\addlinespace
\multicolumn{3}{@{}l}{\emph{Experiment support (goal iv)}} \\
LASSO dataset generator (widths, test battery, $x_0$, ALISTA $W$) & added & \texttt{Benchmarking/data/lasso.py} \\
Experiment-2 metric functions and reference solver & added & \texttt{core/lasso\_metrics.py} \\
Classical ISTA/FISTA baselines on the same instances & added & \texttt{Benchmarking/classical\_baselines.py} \\
Per-iteration relative-loss curves for the learned optimizer & added & \texttt{Benchmarking/plot\_compute\_matched\_rf.py} \\
\bottomrule
\end{tabular}
\end{table}

